\documentclass{article} % Define o tipo de documento (artigo, livro, relatório)

% --- Preâmbulo (onde ficam os pacotes e configurações) ---
\usepackage[utf8]{inputenc} % Acentuação direta
\usepackage{fancyvrb}

\title{Relatório EP1 MAC0352}
\author{Felipe Kelemen, Gustavo Bastos, Lucas Uchiyama, Octavio Carneiro}
\date{\today} % Coloca a data atual automaticamente

% --- Início do Documento ---
\begin{document}

\maketitle % Cria o título com base nas informações acima

\section{Arquitetura do sistema}

O programa utiliza a arquitetura cliente-servidor, juntamente com o protocolo da camada de aplicação DRCP, de nossa autoria, e o protocolo TCP nada camada de transporte, utilizando como base um exemplo de cliente e servidor simples.

O servidor é dividido de forma a separar a manipulação dos recursos (resource manager), a interpretação das mensagens  recebidas (parser), a chamada ao resource manager (handler), e a escuta do recebimento das mensagens (listener).

\section{Protocolo de comunicação}

Foi utilizado como protocolo da camada de aplicação o Distributed Resource Control Protocol (DRCP), criado por nós mesmos.

\subsection{Métodos}

\begin{itemize}
	\item CREATE(valor(string/int))
	      \subitem Caso sucesso: Devolve id do recurso criado(int)
	      \subitem Casos de erro: Limite de recursos (5), Valor inválido (6)

	\item GET(id(int))
	      \subitem Caso sucesso: Devolve valor do recurso(string/int)
	      \subitem Casos de erro: Recurso já reservado por alguém (1), Recurso inexistente (2), Recurso livre porém não reservado (3)

	\item SET(id(int), valor(string/int))
	      \subitem Caso sucesso: Não devolve nada
	      \subitem Casos de erro: Recurso já reservado por alguém (1), Recurso inexistente (2), Recurso livre porém não reservado (3), Valor inválido (6)

	\item RESERVE(id(int))
	      \subitem Caso sucesso: Não devolve nada
	      \subitem Casos de erro: Recurso já reservado por alguém (1), Recurso inexistente (2)

	\item - RELEASE(id(int))
	      \subitem Caso sucesso: Não devolve nada
	      \subitem Casos de erro: Recurso já reservado por alguém (1), Recurso inexistente (2), Recurso livre porém não reservado (3)

	\item - LIST()
	      \subitem Caso sucesso: Devolve contador de ids(int), lista de ids(Resource[])
	      \subitem Casos de erro: Nenhum
\end{itemize}

\subsection{Solicitação e resposta}

A solicitação do cliente para o servidor é realizada passando o nome do método os parâmetros em uma única linha:

\begin{Verbatim}[tabsize=4]
	{Método} {Parâmetros}
\end{Verbatim}

Por exemplo:

\begin{itemize}
	\item CREATE "Meu recurso muito legal"
	\item GET 10
\end{itemize}

Por outro lado, a resposta do servidor para o cliente estará no seguinte formato, também em uma única linha, caso sucesso:

\begin{Verbatim}[tabsize=4]
	{Emoji joinha} {Resultado}
\end{Verbatim}

Caso erro, a resposta será:

\begin{Verbatim}[tabsize=4]
	{Emoji irritado} {Código de erro} {Descrição do erro}
\end{Verbatim}

\subsection{Códigos de erro}

\begin{itemize}
	\item Recurso já reservado por alguém (1)
	\item Recurso inexistente (2)
	\item Recurso livre porém não reservado (3)
	\item Limite de recursos (5)
	\item Método inexistente (7)
\end{itemize}

\subsection{Logging}

O arquivo de logging armazena todas as requisições realizadas pelos clientes com timestamp, ID da conexão, tipo de evento e detalhes da operação. O formato é estruturado como segue:

\begin{Verbatim}[tabsize=4]
	[YYYY-MM-DD HH:MM:SS][#ID_CONEXÃO][EVENT_TYPE] dados
\end{Verbatim}

Exemplos reais de operações bem-sucedidas:

\begin{Verbatim}[tabsize=4]
[2026-04-13 20:17:00][#00001][CONNECT]
[2026-04-13 20:17:00][#00001][REQUEST] CREATE valor_inicial
[2026-04-13 20:17:00][#00001][ANSWER] 0 {OK} 0
[2026-04-13 20:17:00][#00001][DISCONNECT]

[2026-04-13 20:17:00][#00008][REQUEST] RESERVE 7
[2026-04-13 20:17:00][#00008][ANSWER] 0 {OK}

[2026-04-13 20:17:00][#00003][REQUEST] LIST
[2026-04-13 20:17:00][#00003][ANSWER] 0 {OK} 4 0 1 2 3
\end{Verbatim}

Exemplos de operações com erro:

\begin{Verbatim}[tabsize=4]
[2026-04-13 20:17:00][#00006][REQUEST] GET 999999
[2026-04-13 20:17:00][#00006][ANSWER] 2 {ERROR} 2 Recurso inexistente

[2026-04-13 20:17:00][#00005][REQUEST] GET 5
[2026-04-13 20:17:00][#00005][ANSWER] 3 {ERROR} 3 Recurso nao reservado

\begin{Verbatim}[tabsize=4]
	typedef struct{
			int id;
			char* value;
			bool reserved;
			pthread_t* belongs_to;
		} Resource;
\end{Verbatim}

\section{Estratégia de concorrência}

A estratégia de concorrência utilizada foi o uso de semáforos juntamente com threads, permitindo que apenas uma thread tenha acesso aos recursos do servidor por vez.

\section{Tratamento de falhas}

O sistema lida com a falha em cliente tornando todos os recursos reservados por ele como públicos, permitindo que qualquer outro cliente os reserve.

\section{Experimentos realizados}

\section{Limitações e melhorias futuras}



\end{document}

